\chapter{Generative Adversarial Networks}
\label{ch:gans}

\begin{tcolorbox}[feynbox={Sky}{BBg},title={Before and After}]
\textbf{Before this chapter you need:} norms (\S\ref{sec:norms}), minimax and
penalties (\S\ref{sec:lagrange}--\S\ref{sec:minimax}), neural networks
(Chapter~\ref{ch:neuralnets}), cross-entropy, JS divergence, Wasserstein distance
and Kantorovich--Rubinstein duality
(\S\ref{sec:crossentropy}--\S\ref{sec:duality}).\\
\textbf{After it you will understand:} what the two networks optimise, why the
game oscillates rather than converging, what mode collapse looks like in this
project's own measurements, what Lipschitz continuity asks for, and why equal
generator updates across models is emphatically not equal compute.
\end{tcolorbox}

Everything needed to read a GAN objective as a sentence is now in place. This
chapter assembles it and describes the three architectures benchmarked here.

%% ─────────────────────────────────────────────────────────────────────────────
\section{The GAN Idea}
\label{sec:ganidea}

\situation{
A counterfeiter makes fake notes; a teller learns to spot them. Each improvement
on one side forces an improvement on the other. Neither needs a written
definition of ``convincing'' --- the definition is whatever the opponent can
currently detect.
}

\problem{
We want to generate realistic return sequences, and realism has no closed-form
loss. Writing one by hand would mean enumerating the properties that matter,
which is exactly the thing we do not know how to do completely.
}

\workedarith{
\textbf{The insight.} Replace the hand-written loss with a second network trained
to tell real from fake. Its output becomes the generator's training signal.

\medskip
\textbf{What the discriminator says.} Given a sample $x$, it returns
$D(x) \in (0,1)$, read as the probability that $x$ is real.

\medskip
\begin{tabular}{@{}lll@{}}
\toprule
$D(x)$ & Reading & Generator's position \\
\midrule
$0.99$ & confident real & fake was caught easily \\
$0.50$ & cannot tell    & fake is indistinguishable \\
$0.01$ & confident fake & fake was caught easily \\
\bottomrule
\end{tabular}

\medskip
\textbf{The generator's target is $0.5$.} Not 1 --- it is not trying to convince
the discriminator its output is \emph{more} real than real data, only to reach the
point where the discriminator has no information. That is what
$p_G = p_{\text{data}}$ looks like from the discriminator's side.
}

\formaldef{
A \textbf{GAN} (Goodfellow et al.\ 2014) pairs two networks:
\begin{itemize}
  \item \textbf{Generator} $G:\mathcal{Z}\to\mathcal{X}$, mapping noise
        $z \sim p_z$ to a sample $\tilde{x} = G(z)$;
  \item \textbf{Discriminator} $D:\mathcal{X}\to(0,1)$, estimating the probability
        that its input is real.
\end{itemize}
Training is the minimax problem
\[
  \min_G\max_D\; \E_{x\sim p_{\text{data}}}\bigl[\log D(x)\bigr]
  + \E_{z\sim p_z}\bigl[\log\bigl(1 - D(G(z))\bigr)\bigr] .
\]
At the ideal equilibrium $p_G = p_{\text{data}}$ and $D \equiv 1/2$.
}

\analogybreak{
The counterfeiter story implies a winner. There is none: the target state is one
where the discriminator is reduced to guessing. It also implies steady mutual
improvement, and \S\ref{sec:nash} shows the dynamics need not improve
monotonically --- or at all.
}

\whyhere{
All three models here are GANs, chosen because the framework defines the target
distribution implicitly through the adversary rather than requiring an explicit
density (\S\ref{sec:mletraining}). For market returns, whose true density is
unknown and certainly not Gaussian, that is the decisive advantage --- and the
cost is that none of these models can report a likelihood, so evaluation must be
entirely sample-based.
}

%% ─────────────────────────────────────────────────────────────────────────────
\section{The Adversarial Objective and Binary Cross-Entropy}
\label{sec:adversarial}

\situation{
The teller is graded on a simple rule: confident and right earns a good mark,
confident and wrong earns a terrible one, and hedging at $50/50$ earns a mediocre
one whatever the truth.
}

\problem{
The minimax expression of the previous section needs unpacking. Both terms are
instances of one loss --- binary cross-entropy (\S\ref{sec:crossentropy}) --- and
seeing that makes the generator's gradient problem visible.
}

\workedarith{
\textbf{The discriminator's job.} Maximise $\log D(x)$ on real data and
$\log(1-D(\tilde{x}))$ on fakes --- the two halves of binary cross-entropy with
labels 1 and 0.

\medskip
\begin{tabular}{@{}lrrr@{}}
\toprule
$D(\tilde{x})$ on a fake & $\log(1-D)$ & $D$'s score & $G$'s position \\
\midrule
$0.9$  & $\log 0.1 = -2.303$ & poor  & fooling it \\
$0.5$  & $\log 0.5 = -0.693$ & fair  & indistinguishable \\
$0.1$  & $\log 0.9 = -0.105$ & good  & caught \\
$0.01$ & $\log 0.99 = -0.010$ & excellent & caught easily \\
\bottomrule
\end{tabular}

\medskip
\textbf{Now the generator's gradient.} $G$ minimises $\log(1-D(G(z)))$. Its
derivative with respect to $D$ is
\[
  \frac{d}{dD}\log(1-D) = \frac{-1}{1-D} .
\]

\medskip
\begin{tabular}{@{}lrr@{}}
\toprule
$D(\tilde{x})$ & $-1/(1-D)$ & signal strength \\
\midrule
$0.5$  & $-2.00$   & strong \\
$0.1$  & $-1.11$   & moderate \\
$0.01$ & $-1.010$  & weak \\
\bottomrule
\end{tabular}

\medskip
\textbf{The problem.} When the generator is bad --- which is exactly at the start
--- the discriminator confidently returns $D \approx 0$, and the gradient shrinks
toward $-1$ while the \emph{loss} is near zero. The generator learns most slowly
precisely when it has most to learn.

\medskip
\textbf{The standard patch.} Have $G$ maximise $\log D(G(z))$ instead of
minimising $\log(1-D(G(z)))$. Same fixed point, but the derivative $1/D$ is
\emph{large} when $D$ is small: at $D = 0.01$ it is $100$. This is the
``non-saturating'' loss used in practice.
}

\formaldef{
The value function is
\[
  V(D,G) = \E_{x\sim p_{\text{data}}}[\log D(x)] + \E_{z}\bigl[\log(1-D(G(z)))\bigr] .
\]
For a fixed $G$ the optimal discriminator is
\[
  D^{*}(x) = \frac{p_{\text{data}}(x)}{p_{\text{data}}(x) + p_G(x)} ,
\]
and substituting it back gives
\[
  V(D^{*},G) = 2\,\mathrm{JSD}\bigl(p_{\text{data}} \parallel p_G\bigr) - 2\log 2 .
\]
So the original GAN minimises the Jensen--Shannon divergence.
}

\analogybreak{
That last identity is the whole trouble. Section~\ref{sec:jsd} showed JSD is
constant at $\log 2$ whenever the two distributions have disjoint support ---
which is the normal situation for a generator mapping $\R^{100}$ into $\R^{128}$.
A constant objective has zero gradient, so the theoretical target provides no
direction. The non-saturating patch above helps in practice without repairing the
underlying geometry.
}

\whyhere{
This is the theoretical case for replacing JSD with Wasserstein distance, taken up
in \S\ref{sec:wgan}. It is also why TimeGAN, which retains a
cross-entropy-style adversarial term, is more sensitive to its pre-training
phases than the two WGAN-GP models are: its Phase 3 inherits exactly this
gradient behaviour.
}

%% ─────────────────────────────────────────────────────────────────────────────
\section{Nash Equilibrium and Why Training Oscillates}
\label{sec:nash}

\situation{
Two people adjusting a shared thermostat, one preferring cold and one warm. Each
responds sensibly to the current setting. The temperature never settles.
}

\problem{
Ordinary training has a target: minimise the loss, and a falling loss means
progress. In a two-player game there is no single quantity that should fall, so
the usual reading of a training curve does not apply.
}

\workedarith{
\textbf{The equilibrium.} A Nash equilibrium is a pair where neither player
improves by moving alone. For a GAN the ideal is $p_G = p_{\text{data}}$ with
$D \equiv 1/2$, giving
\[
  V = \log\tfrac12 + \log\tfrac12 = -1.386 = -2\log 2 ,
\]
consistent with $\mathrm{JSD} = 0$ in the identity above. \checkmark

\medskip
\textbf{Why it need not be reached.} Section~\ref{sec:minimax} showed the
canonical failure on $L = xy$: simultaneous gradient steps multiply the distance
from equilibrium by $\sqrt{1+\eta^2} > 1$ at every step, for any positive learning
rate. The players spiral outward. Nothing about that argument depends on the
objective being $xy$ specifically --- it is a property of opposed gradients.

\medskip
\textbf{What this means for loss curves.} Consider three scenarios and what the
losses look like:

\medskip
\begin{tabular}{@{}lll@{}}
\toprule
Situation & $D$ loss & $G$ loss \\
\midrule
Healthy training      & oscillating & oscillating \\
Discriminator wins    & near 0      & large \\
Mode collapse         & near 0      & \emph{stable} \\
\bottomrule
\end{tabular}

\medskip
\textbf{Read the last row.} A collapsed generator produces one output, which the
discriminator learns to reject perfectly, so both losses go flat. Flat losses ---
the thing that signals success in ordinary training --- here signal the worst
outcome.
}

\formaldef{
A \textbf{Nash equilibrium} of $\min_G\max_D V(D,G)$ is a pair $(G^{*},D^{*})$
with
\[
  V(D^{*},G^{*}) \leq V(D^{*},G) \ \forall G ,
  \qquad
  V(D^{*},G^{*}) \geq V(D,G^{*}) \ \forall D .
\]
Existence is guaranteed for convex--concave $V$ over compact convex strategy sets.
A GAN's $V$ is neither convex in $G$'s parameters nor concave in $D$'s, so neither
existence nor convergence of simultaneous gradient descent is guaranteed.
}

\analogybreak{
``The GAN converged'' is not a claim these losses can support. Practitioners
routinely stop on sample quality or an external metric instead, and that is not
laziness --- it is the only defensible criterion available when the objective
itself is uninformative.
}

\whyhere{
This is why the pipeline trains for a \emph{fixed gradient-step budget} rather
than to convergence, and why it carries explicit degeneracy guards: a check on the
generated series' mean, standard deviation and ACF(1), and an autoencoder $R^2$
threshold for TimeGAN. All model comparison happens in the external evaluation
framework, never on training loss.
}

%% ─────────────────────────────────────────────────────────────────────────────
\section{Mode Collapse}
\label{sec:modecollapse}

\situation{
A student who has learned one answer that usually earns a pass writes it for every
question. Marks are mediocre but never terrible, and the strategy is stable ---
which is precisely why they never learn anything else.
}

\problem{
The generator is rewarded for fooling the discriminator, not for variety. Nothing
in the objective requires it to cover the data distribution, so collapsing onto a
narrow region can be a locally optimal strategy.
}

\workedarith{
\textbf{This project's case.} During a BOVESPA run, TimeGAN produced a series
with:

\medskip
\begin{tabular}{@{}lrrl@{}}
\toprule
Statistic & Generated & Real & Verdict \\
\midrule
standard deviation & $0.0002$ & $0.0103$ & $51\times$ too small \\
$\mathrm{ACF}(1)$  & $0.831$  & $\approx 0.03$ & wildly autocorrelated \\
\bottomrule
\end{tabular}

\medskip
\textbf{Reading the two numbers together.} A standard deviation $51\times$ too
small means almost no variation --- the generator emits nearly the same value
every day. An $\mathrm{ACF}(1)$ of $0.831$ against a real $0.03$ means consecutive
values are strongly related, so what is being produced is a smooth drifting
\emph{path}, not a return series.

\medskip
\textbf{Why each alone would be ambiguous.} A small standard deviation could in
principle be a calm regime. A high ACF could be a modelling choice. Together they
are unmistakable: a nearly constant, slowly wandering output.

\medskip
\textbf{The guard.} The pipeline checks every generated series against
\[
  |\text{mean}| < 0.01 , \qquad
  0.001 < \text{sd} < 0.15 , \qquad
  |\mathrm{ACF}(1)| < 0.3 ,
\]
and prints a warning on failure. The collapsed run above fails two of the three.
In the later smoke test, TimeGAN produced $\mathrm{ACF}(1) = 0.429$ and $0.443$ ---
still above the $0.3$ threshold, still flagged, but far from the $0.831$ of the
full collapse.

\medskip
\textbf{And the reason metrics alone are insufficient.} Section~\ref{sec:duality}
noted the Wasserstein floor. A constant-output generator scores $W_1$ equal to the
real data's mean absolute deviation --- $0.00770$ for BOVESPA --- which a genuine
but imperfect generator can easily exceed. Without the guards, collapse can win on
individual metrics.
}

\formaldef{
\textbf{Mode collapse} is the failure in which $G$ maps many distinct $z$ to few
distinct outputs, so $p_G$ covers only part of $p_{\text{data}}$'s support. In
the extreme, $G$ is nearly constant.

Contributing causes: the objective rewards fooling $D$ rather than coverage;
$\min_G\max_D$ solved in the wrong order favours a point mass at the discriminator's
current weak spot; and reverse-KL-like objectives are mode-seeking
(\S\ref{sec:kl}).
}

\analogybreak{
Not every low-diversity output is mode collapse. A generator trained on a genuinely
calm period should produce low-variance output, and calling that collapse would be
wrong. What distinguishes the two is the comparison against the real data's own
statistics on the \emph{same} period --- which is why the guard thresholds are
anchored to measured BRICS values, $|{\rm mean}| < 0.001$ and daily volatility
$0.9$--$2.0\%$, rather than chosen abstractly.
}

\whyhere{
Mode collapse is the single most consequential failure mode in this benchmark,
because a collapsed model can score well on several fidelity metrics by accident.
Three independent safeguards address it: the post-generation guard above, the
Hill $\alpha > 20$ rejection of \S\ref{sec:hillfail}, and the composite ranking
design of Chapter~\ref{ch:evaluation}.
}

%% ─────────────────────────────────────────────────────────────────────────────
\section{Lipschitz Continuity}
\label{sec:lipschitz}

\situation{
A hiking trail with a posted maximum gradient. Whatever route you take between two
points, the trail never climbs faster than the posted limit --- so knowing the
horizontal distance between two points bounds how far apart they can be
vertically.
}

\problem{
Kantorovich--Rubinstein duality (\S\ref{sec:duality}) required the supremum to be
taken over $1$-Lipschitz functions, and without that restriction the supremum is
infinite and measures nothing. Making the constraint precise is the last step
before WGAN.
}

\workedarith{
\textbf{A Lipschitz function.} $f(x) = 3x$. For any $a, b$:
\[
  |f(a) - f(b)| = |3a - 3b| = 3|a-b| .
\]
So the ratio $|f(a)-f(b)|/|a-b|$ is exactly $3$ everywhere: $f$ is
$3$-Lipschitz, and $3$ is the smallest such constant.

\medskip
\textbf{Check it on numbers.} $a = 1$, $b = 4$: $|3 - 12| = 9$ and
$3|1-4| = 9$. \checkmark

\medskip
\textbf{A function that is not Lipschitz on $\R$.} $f(x) = x^2$. Take
$b = a + 1$:
\[
  |f(a+1) - f(a)| = |a^2 + 2a + 1 - a^2| = |2a + 1| ,
\]
while $|b - a| = 1$. The ratio is $|2a+1|$, which grows without bound as $a$ does:

\medskip
\begin{tabular}{@{}lrr@{}}
\toprule
$a$ & ratio $|2a+1|$ & \\
\midrule
$1$    & $3$    & \\
$10$   & $21$   & \\
$1000$ & $2001$ & no constant works \\
\bottomrule
\end{tabular}

\medskip
No single $K$ bounds it, so $x^2$ is not Lipschitz on $\R$. It \emph{is} Lipschitz
on any bounded interval --- on $[-5,5]$ the derivative $2x$ is capped at $10$, so
it is $10$-Lipschitz there. Lipschitz continuity is a property of a function
\emph{together with} its domain.

\medskip
\textbf{The derivative connection.} For differentiable $f$, the smallest valid $K$
is $\sup_x |f'(x)|$. For $3x$ that is $3$; for $x^2$ on $\R$ it is unbounded.
In several variables, $|f'|$ becomes $\norm{\nabla f}_2$ --- which is where the
gradient norm enters.
}

\formaldef{
$f$ is \textbf{$K$-Lipschitz} on a domain if
\[
  |f(a) - f(b)| \leq K\,\norm{a - b} \quad\text{for all } a,b
\]
in that domain. The smallest such $K$ is the \textbf{Lipschitz constant}. For
differentiable $f$,
\[
  K = \sup_{x}\norm{\nabla f(x)}_2 .
\]
A $1$-Lipschitz function is one whose gradient norm never exceeds $1$ --- it
cannot amplify distances.
}

\analogybreak{
Lipschitz continuity bounds the \emph{rate} of change, not the size of the
function. A $1$-Lipschitz function can take arbitrarily large values; it just
cannot get there quickly. And the constraint is global: satisfying it on the
training data says nothing about points in between, which is exactly the gap the
gradient penalty's interpolation strategy is designed to cover.
}

\whyhere{
This is the constraint that makes the critic's output a valid Wasserstein
estimate. Because the smallest $K$ equals $\sup\norm{\nabla f}_2$, enforcing
$1$-Lipschitz continuity means controlling the gradient norm --- and that is
exactly what $\norm{\nabla D}_2 = 1$ asks for.
}

%% ─────────────────────────────────────────────────────────────────────────────
\section{WGAN and the Critic}
\label{sec:wgan}

\situation{
Replace the teller who answers ``real or fake'' with an appraiser who assigns a
score on an open scale. The appraiser never says ``certainly fake'', so there is
no confident verdict at which the conversation stops being informative.
}

\problem{
The discriminator's saturation (\S\ref{sec:adversarial}) and JSD's flatness on
disjoint supports (\S\ref{sec:jsd}) are the same defect seen twice. Replacing the
objective with $W_1$ fixes both, and duality makes it computable.
}

\workedarith{
\textbf{What changes.}

\medskip
\begin{tabular}{@{}lll@{}}
\toprule
& Original GAN & WGAN \\
\midrule
Final layer   & sigmoid, output in $(0,1)$ & none, output in $\R$ \\
Name          & discriminator & critic \\
Loss          & binary cross-entropy & $\E[D(x_{\text{real}})] - \E[D(x_{\text{fake}})]$ \\
Constraint    & none & $1$-Lipschitz \\
Divergence    & JSD & $W_1$ \\
\bottomrule
\end{tabular}

\medskip
\textbf{The critic loss, read plainly.} The critic maximises the gap between its
average score on real data and on fake data. By duality (\S\ref{sec:duality}),
that maximised gap \emph{is} $W_1$ --- provided the critic stays $1$-Lipschitz.

\medskip
\textbf{Why the gradient does not vanish.} Recall the two point masses of
\S\ref{sec:jsd}: $p$ at $0$, $q$ at $\theta$. JSD is $\log 2$ for every
$\theta \neq 0$, giving zero gradient. $W_1 = |\theta|$, giving gradient $\pm 1$
--- constant, non-zero, and pointing the right way, no matter how far apart the
distributions start.

\medskip
\textbf{The generator's loss.} $G$ minimises $-\E[D(G(z))]$: push the critic's
score on fakes upward. Unlike the original formulation, this gradient does not
degrade when the critic is confident, because the critic has no saturating output
to be confident with.
}

\formaldef{
\textbf{WGAN} (Arjovsky et al.\ 2017) replaces the discriminator with a
\textbf{critic} $D:\mathcal{X}\to\R$ and optimises
\[
  \min_G\max_{\norm{D}_L\leq 1}\;
  \E_{x\sim p_{\text{data}}}[D(x)] - \E_{z\sim p_z}\bigl[D(G(z))\bigr] .
\]
When the critic is exactly $1$-Lipschitz and fully optimised, the inner maximum
equals $W_1(p_{\text{data}}, p_G)$ by Kantorovich--Rubinstein duality.

The original WGAN enforced the constraint by \textbf{weight clipping}, restricting
every weight to $[-c,c]$. This works but is crude: it biases the critic toward
simple functions and is sensitive to $c$.
}

\analogybreak{
The critic's loss is a genuine $W_1$ estimate only when the critic is both exactly
$1$-Lipschitz and fully optimised. Neither holds in practice: a finite network is a
restricted function class, and the critic is trained for a handful of steps per
generator update. The reported critic loss is therefore a lower bound on an
approximation, and it is not comparable across architectures --- which is one more
reason this project judges models externally rather than on training loss.
}

\whyhere{
QuantGAN and FinGAN both use WGAN-GP. TimeGAN does not: it retains a
cross-entropy-style adversarial term inside its Phase 3, which is one reason its
behaviour differs and why its pre-training phases matter so much more.
}

%% ─────────────────────────────────────────────────────────────────────────────
\section{The Gradient Penalty}
\label{sec:gp}

\situation{
Enforcing the trail's maximum gradient by inspecting every metre is impossible.
Instead, inspect a random sample of points and fine any violation. Enforcement
becomes statistical rather than absolute.
}

\problem{
The Lipschitz constraint must hold everywhere, and there is no way to impose it
exactly on a neural network. Weight clipping over-restricts. What is needed is a
soft constraint (\S\ref{sec:lagrange}) that pushes the gradient norm toward 1
where it matters.
}

\workedarith{
\textbf{The penalty.}
\[
  \mathcal{L}_{\text{GP}} = \lambda\,\E_{\hat{x}}\Bigl[\bigl(\norm{\nabla_{\hat{x}}D(\hat{x})}_2 - 1\bigr)^2\Bigr] ,
  \qquad \lambda = 10 .
\]

\medskip
\textbf{What it costs.} Suppose at some point the gradient norm is:

\medskip
\begin{tabular}{@{}lrr@{}}
\toprule
$\norm{\nabla D}_2$ & $(\norm{\nabla D}_2 - 1)^2$ & penalty at $\lambda = 10$ \\
\midrule
$1.0$ & $0.00$ & $0.0$ \\
$1.5$ & $0.25$ & $2.5$ \\
$0.5$ & $0.25$ & $2.5$ \\
$3.0$ & $4.00$ & $40.0$ \\
\bottomrule
\end{tabular}

\medskip
\textbf{Note it is two-sided.} A norm of $0.5$ is penalised as much as $1.5$. This
is deliberate: the optimal critic under duality \emph{saturates} the constraint,
achieving slope exactly 1 along the transport direction --- as $f(x) = -x$ did in
\S\ref{sec:duality}. A critic with a systematically small gradient is
under-using its budget and gives a weaker signal.

\medskip
\textbf{Where it is evaluated.} Not on real data, and not on fakes, but on random
interpolations:
\[
  \hat{x} = \varepsilon x_{\text{real}} + (1-\varepsilon)x_{\text{fake}} ,
  \qquad \varepsilon \sim U[0,1] .
\]
The straight line between a real and a fake sample is where the transport plan
moves mass, so it is where the Lipschitz constraint actually binds. Enforcing it
on the data alone would leave the space between unconstrained.

\medskip
\textbf{A concrete interpolation.} With $x_{\text{real}} = 0.02$,
$x_{\text{fake}} = -0.01$ and $\varepsilon = 0.3$:
\[
  \hat{x} = 0.3(0.02) + 0.7(-0.01) = 0.006 - 0.007 = -0.001 .
\]
}

\formaldef{
\textbf{WGAN-GP} (Gulrajani et al.\ 2017) trains the critic on
\[
  \E_{\tilde{x}}[D(\tilde{x})] - \E_{x}[D(x)]
  + \lambda\,\E_{\hat{x}}\Bigl[\bigl(\norm{\nabla_{\hat{x}}D(\hat{x})}_2 - 1\bigr)^2\Bigr] ,
\]
with $\hat{x}$ interpolated between real and generated samples and
$\lambda = 10$ standard.

The critic must not use \textbf{batch normalisation}: the penalty is a statement
about the gradient at a single point, and batch normalisation makes each output
depend on the rest of the batch, so the per-point gradient norm is no longer
well defined. Layer normalisation is the usual substitute.
}

\analogybreak{
The penalty enforces the constraint only on the sampled interpolation lines and
only in expectation. Nothing guarantees $1$-Lipschitz behaviour elsewhere, so the
critic is not genuinely $1$-Lipschitz and its output is not exactly $W_1$. This is
a heuristic that works well, not a theorem being applied --- and the distinction
matters when reading a critic loss as though it were a distance.
}

\whyhere{
Both QuantGAN and FinGAN use $\lambda = 10$. FinGAN's critic uses layer
normalisation rather than batch normalisation for exactly the reason above, and
that is a correctness requirement rather than a tuning preference.
}

%% ─────────────────────────────────────────────────────────────────────────────
\section{$n_{\text{critic}}$ and the Compute Asymmetry}
\label{sec:ncritic}

\situation{
Two athletes are to be compared fairly, so each is given the same number of
training sessions. One trains alone; the other brings five sparring partners to
every session. The session counts match. The effort does not.
}

\problem{
Section~\ref{sec:sgd} established that budgets must be specified in gradient
steps. But WGAN-GP takes several critic steps per generator step, so equal
generator budgets do not imply equal work --- and the benchmark must state which
quantity it is holding fixed.
}

\workedarith{
\textbf{Why $n_{\text{critic}} > 1$.} Duality gives $W_1$ only when the critic is
\emph{optimised}. An under-trained critic supplies a biased estimate and therefore
a misleading gradient. Gulrajani et al.\ recommend $n_{\text{critic}} = 5$ as a
practical floor, and this project uses it.

\medskip
\textbf{The arithmetic of the budget.} At $9{,}000$ generator updates:

\medskip
\begin{tabular}{@{}lrrr@{}}
\toprule
Model & Generator steps & $n_{\text{critic}}$ & Critic steps \\
\midrule
TimeGAN  & $9{,}000$ & --- & $9{,}000$ (one per generator step) \\
QuantGAN & $9{,}000$ & $5$ & $45{,}000$ \\
FinGAN   & $9{,}000$ & $5$ & $45{,}000$ \\
\bottomrule
\end{tabular}

\medskip
\textbf{And the wall-clock.} Measured on this project's hardware at $1{,}000$
generator steps:

\medskip
\begin{tabular}{@{}lrr@{}}
\toprule
Model & Seconds & Relative \\
\midrule
TimeGAN  & $19.7$  & $1.0\times$ \\
FinGAN   & $39.3$  & $2.0\times$ \\
QuantGAN & $180.1$ & $9.1\times$ \\
\bottomrule
\end{tabular}

\medskip
\textbf{A $9\times$ spread at identical generator budgets.} The critic ratio
explains part of it and architecture the rest --- QuantGAN's dilated convolutions
over 128 steps are simply more expensive per step than TimeGAN's 24-unit GRU.

\medskip
\textbf{Which is the fair comparison?} Neither, on its own. Equalising generator
updates controls the quantity that defines how much the generator learns, and it
is the choice made here. Equalising wall-clock would answer a different and also
legitimate question --- which model is best per unit of compute --- and would give
TimeGAN roughly nine times as many steps as QuantGAN. The honest course is to fix
one, report the other, and say which was fixed.
}

\formaldef{
WGAN-GP performs $n_{\text{critic}}$ critic updates per generator update. Both
critic and generator use Adam, conventionally with $\beta_1 = 0$ and
$\beta_2 = 0.9$.

$\beta_1 = 0$ disables momentum deliberately. In an adversarial game the useful
gradient direction changes whenever the opponent updates, so accumulated momentum
from previous steps points along a direction that is no longer relevant.
}

\analogybreak{
$n_{\text{critic}} = 5$ is an empirical floor, not a derived optimum. Too few
critic steps and the Wasserstein estimate is biased; too many and compute is
wasted while the generator sees a critic that has over-fitted the current batch of
fakes. No value makes the estimate exact, because the critic is a finite network
trained for finitely many steps (\S\ref{sec:wgan}).
}

\whyhere{
The pipeline asserts that all three models receive equal \emph{generator updates}
and refuses to run otherwise, printing the critic-step counts and TimeGAN's
pre-training steps alongside so the asymmetry is visible rather than hidden.
Wall-clock is recorded per run and reported next to the step budget, because
budget and cost are different claims and the Methods table needs both.
}

%% ─────────────────────────────────────────────────────────────────────────────
\section{TimeGAN}
\label{sec:timegan}

\situation{
Teaching someone to write convincing diary entries. First they learn to summarise
an entry and reconstruct it from the summary. Then they learn to predict tomorrow's
summary from today's. Only once both skills exist do you introduce a critic.
}

\problem{
A GAN operating directly on raw sequences must learn representation and temporal
structure simultaneously through an adversarial signal alone. TimeGAN's answer is
to learn both first, supervised, and let the adversary refine what is already
there.
}

\workedarith{
\textbf{The four components.}

\medskip
\begin{tabular}{@{}lll@{}}
\toprule
Component & Map & Role \\
\midrule
Embedder $e$   & $\mathcal{X}\to\mathcal{H}$ & compress a sequence to latent \\
Recovery $r$   & $\mathcal{H}\to\mathcal{X}$ & decode latent back to data \\
Supervisor $s$ & $\mathcal{H}\to\mathcal{H}$ & predict the next latent step \\
Generator $G$  & $\mathcal{Z}\to\mathcal{H}$ & noise to latent sequence \\
\bottomrule
\end{tabular}

\medskip
\textbf{The phases.} Phase 1 trains the autoencoder to minimise
$\norm{X - r(e(X))}^2$. Phase 2 trains the supervisor on
$\norm{H_{t+1} - s(H_t)}^2$. Phase 3 trains everything jointly against a
discriminator.

\medskip
\textbf{Why Phase 1 must succeed first.} If the autoencoder cannot reconstruct,
the latent space carries no information about the data, and Phase 3's adversarial
game is played entirely in a meaningless space.

The pipeline therefore measures reconstruction $R^2$ and warns below $0.8$. The
history of this measurement in the project is instructive:

\medskip
\begin{tabular}{@{}lrl@{}}
\toprule
Configuration & $R^2$ & \\
\midrule
Sigmoid latent, one batch per epoch & $\approx 0$ & latent stuck near $0.5069$ \\
Tanh latent, full batch loop        & $+0.995$    & healthy \\
\bottomrule
\end{tabular}

\medskip
Both fixes were required: the activation change of \S\ref{sec:activations} and
the batching fix of \S\ref{sec:sgd}.

\medskip
\textbf{Architecture.} GRU backbone, hidden dimension 24, 3 layers, learning rate
$10^{-3}$, batch size 128. Parameter counts: embedder $9{,}744$, recovery
$10{,}825$, generator $11{,}400$, supervisor $7{,}800$, discriminator $10{,}825$.
}

\formaldef{
\textbf{TimeGAN} (Yoon et al.\ 2019) trains in phases, with the joint generator
loss
\[
  \mathcal{L}_G = \mathcal{L}_U + 100\sqrt{\mathcal{L}_S} + 100\,\mathcal{L}_V ,
\]
where $\mathcal{L}_U$ is the unsupervised adversarial term, $\mathcal{L}_S$ the
supervised latent-prediction term, and $\mathcal{L}_V$ a moment-matching term on
mean and standard deviation.

Phases 1 and 2 are \textbf{pre-training} required by the algorithm. This project
budgets them separately --- $3{,}000$ steps each --- and excludes them from the
generator-update parity check of \S\ref{sec:ncritic}, reporting them alongside.
}

\analogybreak{
The latent supervisor is TimeGAN's central idea and its main fragility. It only
helps if the latent space is informative, so the model is unusually sensitive to
the quality of its pre-training --- more so than the two WGAN-GP models, which have
no separate representation stage to get wrong. The $R^2$ guard exists because
failure here is silent: Phase 3 runs perfectly happily on a useless embedding and
reports plausible-looking losses.
}

\whyhere{
TimeGAN's earlier poor showing in this project reflected two implementation
defects rather than the architecture: sigmoid activations throttling the latent
gradient, and a missing batch loop giving it a fraction of the intended training.
Both are fixed, and the current results are the first that can be read as
evidence about the architecture itself.
}

%% ─────────────────────────────────────────────────────────────────────────────
\section{QuantGAN}
\label{sec:quantgan}

\situation{
Rather than reading a sequence step by step, view the whole window at several zoom
levels at once --- one filter seeing adjacent days, another every second day,
another every fourth --- and combine them.
}

\problem{
Recurrent models process sequences serially and their memory decays
(\S\ref{sec:rnn}). Convolutions are parallel and stable but see a narrow window
unless made very deep. Dilation buys reach without depth.
}

\workedarith{
\textbf{The architecture.} Three temporal blocks, each containing \emph{two}
causal convolutions of kernel size 3, with dilations $1, 2, 4$ and channel widths
$[64, 64, 32]$.

\medskip
\textbf{Receptive field, computed properly.} Each convolution with kernel $k$ and
dilation $d$ extends the field by $(k-1)d$. With two convolutions per block:

\medskip
\begin{tabular}{@{}lrrr@{}}
\toprule
Block & Dilation & Extension $2(k-1)d$ & Cumulative \\
\midrule
--- & --- & --- & $1$ \\
0 & $1$ & $4$  & $5$ \\
1 & $2$ & $8$  & $13$ \\
2 & $4$ & $16$ & $29$ \\
\bottomrule
\end{tabular}

\[
  R = 1 + 2(k-1)\sum_l d_l = 1 + 4(1 + 2 + 4) = 29 .
\]

\medskip
\textbf{Twenty-nine days, not eight.} A common shortcut quotes $2^{K}$ for $K$
blocks, giving $8$ here. That formula ignores both the kernel size and the second
convolution in each block, and understates the true field by a factor of nearly
four.

\medskip
\textbf{What 29 days means.} The output at any position depends on roughly six
trading weeks of history --- and on nothing before that. Section~\ref{sec:acf}
measured $\hat\rho(|r|, 50) = +0.05$, still around twice the significance band, so
there is genuine structure beyond the horizon this architecture can represent.
}

\formaldef{
\textbf{QuantGAN} (Wiese et al.\ 2020) uses temporal convolutional networks for
both generator and critic, trained with WGAN-GP.

Configuration here: noise dimension 100, sequence length 128, batch size 64,
learning rate $10^{-4}$ for both networks, $n_{\text{critic}} = 5$,
$\lambda = 10$.
}

\analogybreak{
A large receptive field bounds what the model \emph{can} attend to, not what it
does. The 29-day figure is a ceiling. It is also a hard one: dependence beyond 29
days cannot be represented at all, which sits uncomfortably beside the long-memory
findings of \S\ref{sec:hurst}, where volatility autocorrelation remains
measurable well past lag 50. Extending the field requires a fourth block --- at
dilation 8, adding $32$ to give $61$.
}

\whyhere{
QuantGAN led the composite ranking in this project's smoke test with a
composite rank of $1.50$ against FinGAN's $1.93$ and TimeGAN's $2.57$. Those are
smoke-test numbers at $1{,}000$ generator steps on one market and one seed, and
they are not the benchmark result. The receptive-field limit above is a concrete,
falsifiable prediction about where QuantGAN should underperform in the full run:
on \texttt{hurst\_diff} and long-lag ACF, not on short-horizon clustering.
}

%% ─────────────────────────────────────────────────────────────────────────────
\section{FinGAN}
\label{sec:fingan}

\situation{
Image GANs begin with a coarse grid and upsample repeatedly, adding detail at each
stage. FinGAN applies the same idea to sequences: start from a short noise vector
and expand it to a full 128-step series.
}

\problem{
The benchmark needs a baseline representing the standard convolutional
image-generation approach transferred to time series, deliberately simple, so that
the effect of its inductive bias can be isolated.
}

\workedarith{
\textbf{The generator.} Three \texttt{ConvTranspose1d} blocks with stride 2, base
channels 64. Each doubles the sequence length:
\[
  16 \to 32 \to 64 \to 128 .
\]
This is why the sequence length must be divisible by $2^3 = 8$: at $128$,
$128/8 = 16$ is the starting length. A sequence length of $100$ would fail, since
$100/8 = 12.5$.

\medskip
\textbf{The critic} mirrors this with strided \texttt{Conv1d} layers and
\textbf{layer} normalisation --- not batch normalisation, for the reason given in
\S\ref{sec:gp}: the gradient penalty is a per-point quantity, and batch
normalisation would make each output depend on the rest of the batch.

\medskip
\textbf{Cost.} At $1{,}000$ generator steps FinGAN took $39.3$s against
QuantGAN's $180.1$s --- roughly a fifth of the compute for the same generator
budget, since transposed convolutions over a short latent sequence are far cheaper
than dilated convolutions over 128 steps.
}

\formaldef{
\textbf{FinGAN}, as implemented here, is a transposed-convolution WGAN-GP:
generator of three \texttt{ConvTranspose1d} blocks with stride 2; critic mirroring
with \texttt{Conv1d} and layer normalisation; noise dimension 100, batch size 64,
learning rate $10^{-4}$, $n_{\text{critic}} = 5$, $\lambda = 10$.

Constraint: $\texttt{seq\_len}$ must be divisible by 8.
}

\analogybreak{
The upsampling bias assumes hierarchical structure --- coarse features first, fine
detail as refinement. Images have this; return series plausibly do not. A day's
idiosyncratic move is not a refinement of the month's trend, it is a separate
shock. If FinGAN underperforms QuantGAN, that mismatch is the natural explanation,
and it is a hypothesis the benchmark can test rather than an assumption.
}

\whyhere{
FinGAN is the CNN baseline. Its architecture is deliberately plain --- three
upsampling blocks, no attention, no recurrence --- so that its results speak to the
upsampling inductive bias rather than to architectural sophistication. The
divisibility constraint is a real precondition to check before running on any new
sequence length.
}
